\clearpage\fancyhead[L]{\small PALLAS.}\hypertarget{native-384-lower}{}\label{native-384-lower}
\begin{center}\textbf{PALLAS.}\par\vspace{4pt}\hrule\vspace{4pt}\fontsize{10}{12}\selectfont Göttingische gelehrte Anzeigen. Stück 199. Seite 1985\;..\;1990. 1814. Dec. 12.\par\vspace{4pt}\hrule\vspace{4pt}\end{center}

\par\indent Die schon öfters in unsern Blättern erwähnte Arbeit des Herrn Prof. \textsc{Gauss} über die
verwickelte Theorie der Bewegungen der Pallas hat durch die Beobachtungen der im October d. J.
eingetretenen Opposition dieses Planeten — der \emph{zehnten} seit seiner Entdeckung — abermals
eine schöne Bestätigung erhalten. Wir theilen hier diese Beobachtungen und die \emph{zunächst}
daraus abgeleiteten Resultate mit, indem wir uns die weitern Folgerungen in Beziehung auf die
Theorie für eine andere Gelegenheit vorbehalten. Da bei der vor länger als einem Jahre durch
Herrn \textsc{Nicolai}, Adjunct der Seeberger Sternwarte, geführten Vorausberechnung des Laufs
für die diesmalige Sichtbarkeit blos die osculirenden Elemente der neunten Opposition zum
Grunde gelegt waren (was allerdings zur Auffindung vollkommen hinreichend sein musste), so
hatte Hr. Prof. \textsc{Gauss} es für interessant gehalten, dass die Verbesserung dieser
Elemente durch die fortwährende Einwirkung des Jupiter von der neunten bis zur zehnten
Opposition noch vor der Wiedererscheinung des Planeten bestimmt wurde, und diese nicht
unbedeutende Arbeit unter seiner Leitung schon zu Anfang dieses Jahres durch Hrn.
\textsc{Möbius} ausführen lassen, welcher sich damals bei uns mit ausgezeichnetem Eifer den
mathematischen und astronomischen Studien widmete. Von den Resultaten dieser Rechnung, die man
ausführlicher auch schon in dem so eben erschienenen astronomischen Jahrbuche für 1817
abgedruckt findet, setzen wir hier nur, zur Vergleichung mit dem, was die wirklichen
Beobachtungen gegeben haben, die Vorausbestimmung der Opposition selbst her.

\begin{center}1814 October 25. $12^{\mathrm u}\;39^{\mathrm m}\;50^{\mathrm s}5$ Mittlere Zeit in Göttingen\end{center}
\begin{center}\fontsize{10}{12}\selectfont\setlength{\tabcolsep}{3pt}\renewcommand{\arraystretch}{1.15}\begin{tabular}{lcr}
Wahre Länge & $.\;.\;.\;.\;.\;.\;.\;.\;.\;.\;.\;.\;.$ & $31^\circ\;58'\;28''4$\\
Heliocentrische Breite & $.\;.\;.\;.\;.\;.\;.\;.\;.\;.$ & $23\;38\;32.7\;\mathrm{S.}$\\
\end{tabular}\end{center}

\SourceLeaf{385}\Needspace{10\baselineskip}\par\indent Auf der hiesigen Sternwarte wurde folgende Beobachtung mit dem Kreismicrometer gemacht.

\begin{center}\fontsize{9}{11}\selectfont\setlength{\tabcolsep}{3pt}\renewcommand{\arraystretch}{1.15}\begin{tabular}{c|ccc|ccc|ccc}
1814 & \multicolumn{3}{c|}{\shortstack{Mittlere Zeit\\in Göttingen}} & \multicolumn{3}{c|}{\shortstack{Scheinbare gerade\\Aufsteigung}} & \multicolumn{3}{c}{\shortstack{Scheinbare\\Abweichung}} \\
\hline
Sept.\ 16 & $12^{\mathrm u}$ & $1^{\mathrm m}$ & $1^{\mathrm s}$ & $46^{\circ}$ & $16'$ & $23''4$ & $11^{\circ}$ & $4'$ & $20''9$ S. \\
\end{tabular}\end{center}

\par\indent Da sich hiedurch die Richtigkeit der Ephemeride bestätigte, und nicht zu zweifeln war, dass
auf andern Sternwarten der Planet mit festen Instrumenten würde beobachtet werden können, so
wäre die weitere Fortsetzung der Kreismikrometerbeobachtungen ohne Zweck gewesen. Uebrigens
gründete sich obige Beobachtung auf zahlreiche Vergleichungen, und war im Tagebuche als sehr
gut bezeichnet.

\Needspace{19\baselineskip}\par\indent Auf der \emph{Seeberger} Sternwarte beschränkte man sich mit Recht auf die Beobachtungen am
Mittagsfernrohr, da der Planet zu lichtschwach war, um von dem \textsc{Dollond}schen Quadranten
zuverlässige Declinationen erwarten zu können. Die von Hrn. von \textsc{Lindenau} und Hrn.
\textsc{Nicolai} beobachteten geraden Aufsteigungen sind folgende:

\begin{center}\fontsize{9}{11}\selectfont\setlength{\tabcolsep}{3pt}\renewcommand{\arraystretch}{1.15}\begin{tabular}{c|ccc|ccc}
1814 & \multicolumn{3}{c|}{\shortstack{Mittlere Zeit\\auf Seberg}} & \multicolumn{3}{c}{\shortstack{Scheinbare gerade\\Aufsteigung}} \\
\hline
Sept.\ 17  & $15^{\mathrm u}$ & $19^{\mathrm m}$ & $41^{\mathrm s}5$ & $46^{\circ}$ & $18'$ & $26''2$ \\
18         & 15 & 15 & 52.0 & 46 & 20 & 2.8 \\
19         & 15 & 12 & 1.0  & 46 & 21 & 17.0 \\
20         & 15 & 8  & 8.0  & 46 & 22 & 0.7 \\
Octob.\ 12 & 13 & 36 & 19.2 & 45 & 2  & 4.5 \\
13         & 13 & 31 & 52.0 & 44 & 54 & 14.6 \\
15         & 13 & 22 & 53.1 & 44 & 37 & 25.8 \\
21         & 12 & 55 & 28.0 & 43 & 39 & 51.6 \\
28         & 12 & 22 & 45.3 & 42 & 21 & 48.5 \\
\end{tabular}\end{center}

\Needspace{14\baselineskip}\par\indent Die von Hrn. \textsc{Schumacher} auf der \emph{Mannheimer} Sternwarte mit dem achtfüssigen
Mauerquadranten angestellten Beobachtungen wurden vollständig im Original eingesandt, und
hier von Hrn. \textsc{Encke} auf das sorgfältigste reducirt. Es ergab sich daraus folgendes:

\begin{center}\fontsize{9}{11}\selectfont\setlength{\tabcolsep}{3pt}\renewcommand{\arraystretch}{1.15}\begin{tabular}{c|ccc|ccc|ccc}
1814 & \multicolumn{3}{c|}{\shortstack{Mittlere Zeit\\in Mannheim}} & \multicolumn{3}{c|}{\shortstack{Scheinbare gerade\\Aufsteigung}} & \multicolumn{3}{c}{\shortstack{Scheinbare\\Abweichung}} \\
\hline
Octob.\ 12 & $13^{\mathrm u}$ & $36^{\mathrm m}$ & $17^{\mathrm s}4$ & $45^{\circ}$ & $2'$  & $0''3$  & $19^{\circ}$ & $12'$ & $38''4$ S. \\
15         & 13 & 22 & 51.5 & 44 & 37 & 23.7 & 20 & 6  & 54.6 \\
21         & 12 & 55 & 27.0 & 43 & 39 & 57.5 & 21 & 49 & 52.1 \\
22         & 12 & 50 & 48.3 & 43 & 29 & 15.3 & 22 & 6  & 3.8  \\
\end{tabular}\end{center}

\par\indent Bei den beiden letzten Beobachtungen war der Planet so lichtschwach, dass er sich nur mit
Mühe erkennen liess.

\Needspace{18\baselineskip}\par\indent Auf der neuen \emph{Königsberger} Sternwarte (von deren grosser Thätigkeit unsre Blätter
schon mehrere Proben gegeben haben), machte Hr. Prof. \textsc{Bessel} folgende
Beobachtungen:

\begin{center}\fontsize{9}{11}\selectfont\setlength{\tabcolsep}{3pt}\renewcommand{\arraystretch}{1.15}\begin{tabular}{c|ccc|ccc|ccc}
1814 & \multicolumn{3}{c|}{\shortstack{Mittlere Zeit\\in Königsberg}} & \multicolumn{3}{c|}{\shortstack{Scheinbare gerade\\Aufsteigung}} & \multicolumn{3}{c}{\shortstack{Scheinbare\\Abweichung}} \\
\hline
Octob.\ 14 & $13^{\mathrm u}$ & $27^{\mathrm m}$ & $30^{\mathrm s}6$ & $44^{\circ}$ & $46'$ & $15''3$ & $19^{\circ}$ & $48'$ & $30''8$ S. \\
15         & 13 & 23 & 0.4  & 44 & 37 & 38.6 & 20 & 6  & 20.0 \\
18         & 13 & 9  & 22.9 & 44 & 10 & 8.3  & 20 & 58 & 58.9 \\
19         & 13 & 4  & 48.3 & 44 & 0  & 25.9 & 21 & 15 & 52.6 \\
Nov.\ 1    & 12 & 3  & 57.3 & 41 & 33 & 59.9 & 24 & 31 & 29.8 \\
2          & 11 & 59 & 12.6 & 41 & 21 & 46.2 & 24 & 44 & 9.7  \\
3          & 11 & 54 & 27.4 & 41 & 9  & 24.1 & 24 & 56 & $10\;6$ \\
4          & 11 & 49 & 42.0 & 40 & 56 & 59.6 & 25 & 7  & 41.8 \\
\end{tabular}\end{center}

\SourceLeaf{386}\par\indent Auch dieser vortreffliche Beobachter beklagt sich über die Schwierigkeiten, die die Lichtschwäche des Planeten den Beobachtungen entgegenstellte.

\Needspace{31\baselineskip}\par\indent Die Resultate der Vergleichung dieser sämmtlichen Beobachtungen mit den zuletzt gefundenen Elementen gibt folgende Uebersicht (wobei nur zu bemerken ist, dass die angewandte Epoche der mittleren Länge einige Secunden abweicht, vermöge eines dem Hrn.\ Prof.\ \textsc{Gauss} eigenthümlichen und an einem andern Orte umständlicher zu rechtfertigenden Verfahrens):

\begin{center}
\fontsize{10}{12}\selectfont\setlength{\tabcolsep}{5pt}\renewcommand{\arraystretch}{1.15}\begin{tabular}{l|cc|l}
 & \multicolumn{2}{c|}{Unterschied} & \\
1814 & Ger.\ Aufst. & Abweich. & Beobachter \\
\hline
Sept.\ 16 & $-9''2$ & $+5''9$ & Gauss \\
17 & $-1.6$ & & von Lindenau \\
18 & $-6.2$ & & von Lindenau \\
19 & $-10.7$ & & von Lindenau \\
20 & $-7.5$ & & von Lindenau \\
Octob.\ 12 & $-4.9$ & & von Lindenau \\
12 & $-3.6$ & $+2.2$ & Schumacher \\
13 & $-8.8$ & & von Lindenau \\
14 & $-9.4$ & $+12.0$ & Bessel \\
15 & $-5.2$ & $+6.7$ & Bessel \\
15 & $-6.5$ & & von Lindenau \\
15 & $-7.7$ & $+5.7$ & Schumacher \\
18 & $-2.7$ & $+13.2$ & Bessel \\
19 & $-4.0$ & $+2.4$ & Bessel \\
21 & $-2.2$ & & von Lindenau \\
21 & $-12.0$ & $+4.5$ & Schumacher \\
22 & $-0.4$ & $-1.3$ & Schumacher \\
28 & $+1.8$ & & von Lindenau \\
Nov.\ 1 & $-0.4$ & $+8.5$ & Bessel \\
2 & $-5.1$ & $+10.6$ & Bessel \\
3 & $-5.9$ & $-2.8$ & Bessel \\
4 & $-7.9$ & $-21.9$ & Bessel \\
\end{tabular}
\end{center}

\Needspace{10\baselineskip}\par\indent Das aus diesen Beobachtungen gezogene Resultat für die \emph{zehnte Opposition} selbst ist endlich folgendes:

\begin{center}1814 Octob. 25. $12^{\mathrm u}\;33^{\mathrm m}\;22^{\mathrm s}$ Mittlere Zeit in Göttingen\end{center}
\begin{center}\fontsize{10}{12}\selectfont\setlength{\tabcolsep}{3pt}\renewcommand{\arraystretch}{1.15}\begin{tabular}{lcr}
Wahre Länge & $.\;.\;.\;.\;.\;.\;.\;.\;.\;.\;.\;.$ & $31^\circ\;58'\;11''3$\\
Geocentrische Breite & $.\;.\;.\;.\;.\;.\;.\;.\;.\;.$ & $37\;20\;53.2\;\mathrm{Südl.}$\\
\end{tabular}\end{center}

\Needspace{21\baselineskip}\par\indent Die scharfe Vergleichung dieses Resultats mit den zuletzt gefundenen und durch die Jupiterstörungen gehörig berichtigten Elementen gab nur so kleine Unterschiede, dass die Grundelemente darnach keiner weitern Verbesserung bedürfen. Wir setzen also nur noch die Uebersicht der Uebereinstimmung sämmtlicher bisher beobachteten Oppositionen mit der Theorie des Hrn.\ \textsc{Gauss} hierher.

\begin{center}
\fontsize{10}{12}\selectfont\setlength{\tabcolsep}{5pt}\renewcommand{\arraystretch}{1.15}\begin{tabular}{ll|cc}
 & & \multicolumn{2}{c}{Unterschied} \\
Oppos. & & Mittl.\ Länge & Helioc.\ Breite \\
\hline
I & 1803 & $-1''3$ & $-5''3$ \\
II & 1804 & $-12.9$ & $-11.9$ \\
III & 1805 & $-20.4$ & $-0.4$ \\
IV & 1807 & $+2.5$ & $-6.8$ \\
V & 1808 & $+23.5$ & $+9.0$ \\
VI & 1809 & $+21.5$ & $-3.1$ \\
VII & 1811 & $+14.1$ & $-7.2$ \\
VIII & 1812 & $-16.1$ & $+4.0$ \\
IX & 1813 & $-10.9$ & $-0.0$ \\
X & 1814 & $-14.0$ & $+2.4$ \\
\end{tabular}
\end{center}
